\scanpage{94}
\mspage{\handwritten{1'}}

any program of homological, homotopic\handwritten{\unsure{a}}l or topological algebra.

\margin{\handwritten{2.}}The motivation for this letter was the latter topic however. \deleted{\unsure{Br}}\unsure{Ronnie} Brown and his friends are competent algebraists and apparently strongly motivated for investing energy in foundational work, on the other hand they visibly are lacking the necessary scope of vision which geometry alone provides. They seem to me kind of\unsure{:}isolated, partly due I guess to the disrepute I mentioned before \unsure{o-|-} I suggested to try and have contact with people such as yourself, Larry Breen, Illusie and others, who have the geometric insight and \deleted{\unsure{what}}w\handwritten{ho} moreover, \deleted{I believe} may not think themselves to ggod for indulging in occasional reflection o\handwritten{n} foundational matters\handwritten{,} and in the process help others do the work which should be done.

At first sight it had \deleted{\unsure{appea}}\unsure{seemed|appead} to me that the Bangor group had indeed come to work out (quite independently) one basic intuition of the program I had envisioned in those letters to Larry Breen - namely that the study of $n$-truncated homotopy types (of semisimplicial sets, or of topological spaces) was essentially equivalent to the study of so-called $n$-groupoids (where $n$ is any natural integer). This is \unsure{ex$\hat{\mathrm{p}}$ected|expected} to be \margin{\handwritten{\unsure{``}}}achieved by associating to any space (say) $\mathrm{X}$ it's``fundamental $n$-groupoid'' $\handwritten{\mathbb{\mathord{\rlap{\kern.23em\vrule height .66em width .035em}\mathrm{\Pi}}}}_{n}(\mathrm{X})$, generalizing the familiar Poincaré fundamental groupoid for $n=1$. \margin{\handwritten{\unsure{'}}}The o\unsure{B|b}vious idea is that 0-objects of $\handwritten{\mathbb{\mathord{\rlap{\kern.23em\vrule height .66em width .035em}\mathrm{\Pi}}}}_{n}(\mathrm{X})$ should be the points of $\mathrm{X}$, 1-objects should be ``homotopies'' or paths between points, 2-objects should be homotopies between 1-objects, etc.\unsure{-} This $\handwritten{\mathbb{\mathord{\rlap{\kern.23em\vrule height .66em width .035em}\mathrm{\Pi}}}}_{n}(\mathrm{X})$ should embody the $n$-truncated homotopy type of $\mathrm{X}$, in much the same way as for $n=1$ the usual fundamental groupoid embodies the 1-truncated homotopy type. For two spaces $\mathrm{X},\mathrm{Y}$, the set of homotopy-classes of maps $\mathrm{X}\handwritten{\to}\mathrm{Y}$ (more correctly, for general $\mathrm{X},\mathrm{Y}$, the maps of $\mathrm{X}$ into $\mathrm{Y}$ in the homotopy category) should correspond to $n$-equivalence classes of $n$-functors from $\handwritten{\mathbb{\mathord{\rlap{\kern.23em\vrule height .66em width .035em}\mathrm{\Pi}}}}_{n}(\mathrm{X})$ to $\handwritten{\unsure{\mathbb{\mathord{\rlap{\kern.23em\vrule height .66em width .035em}\mathrm{\Pi}}}}}_{n}(\mathrm{Y})$ - etc. There are very stro\unsure{-}ng suggestions for a nice formalism, including a notion of geometric realization of an $n$-groupoid, which should imply that any $n$-groupoid is $n$-equivalent to a $\handwritten{\mathbb{\mathord{\rlap{\kern.23em\vrule height .66em width .035em}\mathrm{\Pi}}}}_{n}(\mathrm{X})$. Moreover, when the notion of $n$-groupoid (or more generally of an $n$-category) is relativized over an arbitrary topos to the notion of an $n$-gerb (or more generally, an $n$-stack), these become the \deleted{\unsure{main}} natural ``coefficients'' for a formalism of non commutative cohomol\unsure{o}gical algebra, in the spirit of Giraud's thesis.

But all this kind of thing\unsure{s} for the time being is pure heuristics - I never so far sat down to try to make explicit at least a definition of $n$-categories and $n$-groupoids, of $n$-functors between these etc. When I got the Bangor reprints I at once had the feeling that this kind of \handwritten{work}
% The extra upright in each pen Pi is drawn explicitly for the supplied math fonts.
% At Pi_n(Y) the top stroke is lost in this copy.
% archive sheet number: 93
